% ============================================================
% CHAPITRE 6
% ============================================================
\chapter{Sprint 5 : Intelligence artificielle}

\section{Introduction}

Pour clôturer la phase de programmation, le cinquième jalon apporte une dimension futuriste au projet : l'intégration d'un robot conversationnel (chatbot). À l'heure où les clients posent les mêmes questions rébarbatives ("Faites-vous des menus sans porc ?", "Livrez-vous le dimanche ?"), ce module IA vient soulager les gérants d'Assiette Gala. Nous verrons ici comment la technique RAG (Retrieval-Augmented Generation) permet d'empêcher l'intelligence artificielle de mentir sur la carte.

% --------------------------------------------------
\section{Backlog du sprint}

\begin{table}[H]
    \centering
    \caption{Backlog du sprint 5}
    \label{tab:backlog_sprint5}
    \renewcommand{\arraystretch}{1.3}
    \begin{tabular}{|c|p{8cm}|c|}
        \hline
        \textbf{ID} & \textbf{User Story} & \textbf{Points} \\
        \hline
        US-25 & En tant que visiteur, je souhaite poser des questions au chatbot afin d'obtenir des informations sur les plats et services. & 13 \\
        \hline
        US-26 & En tant que client, je souhaite recevoir des recommandations personnalisées. & 8 \\
        \hline
        US-27 & En tant qu'administrateur, je souhaite actualiser la base de connaissances du chatbot. & 3 \\
        \hline
    \end{tabular}
\end{table}

% --------------------------------------------------
\section{Conception}

\subsection{Diagramme de cas d'utilisation}

\begin{figure}[H]
    \centering
    \includegraphics[width=0.8\textwidth]{uc_sprint5.png}
    \caption{Diagramme de cas d'utilisation du sprint 5}
    \label{fig:uc_sprint5}
\end{figure}

\subsection{Description textuelle des cas d'utilisation}

\subsubsection{Cas d'utilisation : Interagir avec le chatbot}

\begin{table}[H]
    \centering
    \caption{Description textuelle --- Interagir avec le chatbot}
    \label{tab:desc_chatbot}
    \renewcommand{\arraystretch}{1.3}
    \begin{tabular}{|p{4cm}|p{8.5cm}|}
        \hline
        \textbf{Cas d'utilisation} & Interagir avec le chatbot \\
        \hline
        \textbf{Acteur principal} & Visiteur / Client \\
        \hline
        \textbf{Précondition} & L'utilisateur accède à l'interface du chatbot \\
        \hline
        \textbf{Postcondition} & L'utilisateur reçoit une réponse pertinente à sa question \\
        \hline
        \textbf{Scénario principal} &
        \begin{enumerate}
            \item L'internaute déplie la bulle de discussion en bas de l'écran
            \item Il tape une question liée à un événement ou un plat
            \item Le système FastAPI convertit la phrase en coordonnées mathématiques (Vecteurs)
            \item Extraction des recettes qui parlent du même sujet dans ChromaDB
            \item Injection de ces ingrédients secrets dans le cerveau de Gemini 2.5
            \item L'IA rédige une phrase complète et polie
            \item Apparition de la réponse dans la fenêtre de l'utilisateur
        \end{enumerate} \\
        \hline
        \textbf{Scénario alternatif} &
        \begin{itemize}
            \item[4a.] Base de données vide : Gemini avoue son ignorance
            \item[6a.] Sujet politique ou hors-traiteur : esquive programmée de l'IA
        \end{itemize} \\
        \hline
    \end{tabular}
\end{table}

\subsubsection{Cas d'utilisation : Recevoir des recommandations}

\begin{table}[H]
    \centering
    \caption{Description textuelle --- Recommandations personnalisées}
    \label{tab:desc_recommandation}
    \renewcommand{\arraystretch}{1.3}
    \begin{tabular}{|p{4cm}|p{8.5cm}|}
        \hline
        \textbf{Cas d'utilisation} & Recevoir des recommandations personnalisées \\
        \hline
        \textbf{Acteur principal} & Client \\
        \hline
        \textbf{Précondition} & Le client est authentifié \\
        \hline
        \textbf{Postcondition} & Une liste de plats recommandés est affichée \\
        \hline
        \textbf{Scénario principal} &
        \begin{enumerate}
            \item Analyse nocturne du fichier client
            \item Extraction de ses habitudes (ex: passion pour les fruits de mer)
            \item Croisement vectoriel pour dénicher des salades ou plats similaires
            \item Poussée de ces nouvelles tuiles culinaires sur son écran d'accueil
        \end{enumerate} \\
        \hline
    \end{tabular}
\end{table}

\subsection{Architecture du pipeline RAG}

Le pipeline RAG (\textit{Retrieval-Augmented Generation}) constitue le mécanisme central du chatbot. Il combine la recherche d'information dans une base de connaissances vectorielle avec la capacité de génération d'un modèle de langage. La figure~\ref{fig:pipeline_rag} illustre ce pipeline.

\begin{figure}[!htbp]
    \centering
    \includegraphics[width=0.95\textwidth]{pipeline_rag.png}
    \caption{Architecture du pipeline RAG du chatbot}
    \label{fig:pipeline_rag}
\end{figure}

Le pipeline se décompose en deux phases :

\subsubsection{Phase d'apprentissage (hors ligne)}

\begin{enumerate}
    \item \textbf{Extraction} : Le script Python aspire le menu complet depuis PostgreSQL (plats, tarifs, traces d'arachides).
    \item \textbf{Formatage textuel} : L'API rédige une mini-histoire pour chaque produit de type : <<~Le gratin de crevettes coûte 25 TND, et contient des crustacés~>>.
    \item \textbf{Calcul matriciel} : HuggingFace (MiniLM) compresse ces phrases en tableaux de 384 dimensions.
    \item \textbf{Classement} : ChromaDB gèle ces points spatiaux sur le disque dur.
\end{enumerate}

\subsubsection{Phase de conversation (en ligne)}

\begin{enumerate}
    \item \textbf{Interception} : Le bot lit : "Avez-vous des plats sans poissons ?".
    \item \textbf{Projection} : Cette question devient à son tour un vecteur de 384 cases.
    \item \textbf{Cosinus mathématique} : Le moteur confronte ce vecteur à tous les plats enregistrés (comparaison géométrique d'angles).
    \item \textbf{Tamisage} : On isole le top 5 des recettes végétariennes ou carnivores trouvées.
    \item \textbf{Orchestration LangChain} : Le contexte "Voici 5 plats sans poisson" est glissé invisiblement dans The Prompt.
    \item \textbf{Synthèse} : Gemini 2.5 Flash articule une réponse humaine ("Oui, je vous conseille notre Cordon Bleu...").
\end{enumerate}

\subsection{Diagrammes de séquence}

\subsubsection{Séquence : Question au chatbot}

\begin{figure}[!htbp]
    \centering
    \includegraphics[width=0.9\textwidth]{seq_chatbot.png}
    \caption{Diagramme de séquence --- Question au chatbot}
    \label{fig:seq_chatbot}
\end{figure}

\subsubsection{Séquence : Actualisation de la base de connaissances}

\begin{figure}[!htbp]
    \centering
    \includegraphics[width=0.9\textwidth]{seq_refresh_knowledge.png}
    \caption{Diagramme de séquence --- Actualisation de la base de connaissances}
    \label{fig:seq_refresh_knowledge}
\end{figure}

\subsection{Diagramme de composants}

\begin{figure}[!htbp]
    \centering
    \includegraphics[width=0.92\textwidth,height=0.72\textheight,keepaspectratio]{class_sprint5.png}
    \caption{Diagramme de composants du service d'intelligence artificielle}
    \label{fig:class_sprint5}
\end{figure}

Le service IA est organisé en modules Python :
\begin{itemize}
    \item \texttt{main.py} : point d'entrée FastAPI, déclaration des routes et cycle de vie de l'application
    \item \texttt{services/chatbot.py} : orchestration du pipeline RAG (classification d'intention, indexation ChromaDB, recherche vectorielle, génération LLM)
    \item \texttt{services/recommender.py} : moteur de recommandation hybride (filtrage collaboratif + contenu)
    \item \texttt{services/database.py} : chargement des données depuis les bases PostgreSQL (pools de connexions)
    \item \texttt{routers/chat.py} : endpoints REST du chatbot avec authentification JWT
    \item \texttt{routers/recommendations.py} : endpoints REST des recommandations
    \item \texttt{models/schemas.py} : modèles Pydantic de validation des requêtes et réponses
    \item \texttt{config.py} : configuration (clés API, URLs des bases de données)
\end{itemize}

\subsection{Diagramme d'états-transitions}

Le diagramme d'états-transitions (figure~\ref{fig:state_sprint5}) modélise le pipeline de traitement d'un message par le chatbot, depuis la réception de la question jusqu'à l'affichage de la réponse.

\begin{figure}[!htbp]
    \centering
    \includegraphics[width=0.75\textwidth]{state_sprint5_chatbot.png}
    \caption{Diagramme d'états-transitions --- Pipeline du chatbot RAG}
    \label{fig:state_sprint5}
\end{figure}

% --------------------------------------------------
\section{Réalisation}

\subsection{Points d'entrée de l'API}

\begin{table}[H]
    \centering
    \caption{Endpoints REST du service IA}
    \label{tab:endpoints_ai}
    \renewcommand{\arraystretch}{1.2}
    \begin{tabular}{|l|p{4.5cm}|p{5.5cm}|}
        \hline
        \textbf{Méthode} & \textbf{Route} & \textbf{Description} \\
        \hline
        POST & /api/chat/message & Envoi d'une question au chatbot \\
        \hline
        GET & /api/chat/history/:userId & Historique de conversation \\
        \hline
        DELETE & /api/chat/history/:userId & Effacer l'historique \\
        \hline
        POST & /api/chat/refresh-knowledge & Actualisation de la base de connaissances (administrateur) \\
        \hline
        GET & /api/chat/metrics & Métriques de qualité de la recherche vectorielle (administrateur) \\
        \hline
        GET & /api/recommendations/plats & Recommandations personnalisées \\
        \hline
        POST & /api/recommendations/event-menu & Suggestions de menu pour événements \\
        \hline
        GET & / & Vérification de l'état du service \\
        \hline
    \end{tabular}
\end{table}

\subsection{Choix technologiques}

\subsubsection{Google Gemini 2.5 Flash}

Le modèle de langage retenu est \textbf{Google Gemini 2.5 Flash}. Ce choix est motivé par :
\begin{itemize}
    \item \textbf{Performance} : temps de réponse rapide, adapté à un usage conversationnel en temps réel
    \item \textbf{Coût} : modèle économique avec un quota gratuit généreux pour le développement
    \item \textbf{Qualité} : capacité de compréhension et de génération en langue française
    \item \textbf{API accessible} : intégration simplifiée via LangChain
\end{itemize}

\subsubsection{Modèle d'embedding : all-MiniLM-L6-v2}

Le modèle d'embedding \texttt{sentence-transformers/all-MiniLM-L6-v2} est choisi pour la vectorisation des documents et des questions. Ses caractéristiques~:
\begin{itemize}
    \item Vecteurs de dimension 384 (compact et efficace)
    \item Entraîné sur plus d'un milliard de paires de phrases
    \item Exécution locale sans appel réseau (confidentialité des données)
    \item Temps d'inférence rapide sur CPU
\end{itemize}

\subsubsection{Base de données vectorielle ChromaDB}

Le stockage des vecteurs utilise \textbf{ChromaDB}, une base de données vectorielle open-source avec persistance sur disque. Ce choix se justifie par :
\begin{itemize}
    \item La persistance native des données (les embeddings survivent aux redémarrages)
    \item Le support des métadonnées pour le filtrage (par type de document, catégorie, intention)
    \item L'indexation HNSW intégrée pour une recherche efficace en haute dimension
    \item L'absence de service externe (mode embarqué, pas de serveur séparé)
    \item La compatibilité avec LangChain via le package \texttt{langchain-chroma}
\end{itemize}

\subsection{Algorithme du pipeline RAG}

Le fonctionnement du pipeline RAG se décompose en deux phases principales :

\textbf{Phase d'indexation :}
\begin{enumerate}
    \item Chargement des données depuis la base de données PostgreSQL (plats, catégories, prix, descriptions).
    \item Transformation de chaque plat en un document textuel structuré contenant l'ensemble de ses attributs.
    \item Génération des plongements vectoriels (embeddings) via le modèle \textit{all-MiniLM-L6-v2}.
    \item Stockage des documents et de leurs vecteurs dans le magasin vectoriel ChromaDB.
\end{enumerate}

\textbf{Phase de requête :}
\begin{enumerate}
    \item Vectorisation de la question de l'utilisateur avec le même modèle d'embeddings.
    \item Recherche des 5 documents les plus similaires par similarité cosinus dans ChromaDB.
    \item Construction du prompt contextuel avec les documents récupérés, la question et les instructions système.
    \item Envoi du prompt au modèle Gemini 2.5 Flash pour génération de la réponse.
    \item Retour de la réponse à l'utilisateur via le widget de chat.
\end{enumerate}



\subsection{Configuration du prompt système}

Le chatbot est paramétré avec un prompt système qui définit son comportement :

\begin{itemize}
    \item Il se présente comme l'assistant de <<~Assiette Gala Sfaxienne~>>
    \item Il répond exclusivement en français
    \item Il s'appuie sur le contexte fourni pour formuler ses réponses
    \item Il oriente les questions hors périmètre vers les canaux de contact du traiteur
    \item Il adopte un ton professionnel, courtois et informatif
\end{itemize}

\subsection{Interfaces réalisées}

\subsubsection{Widget de chat}

\begin{figure}[H]
    \centering
    \includegraphics[width=0.5\textwidth]{ui_chatbot.png}
    \caption{Widget de chat du chatbot IA}
    \label{fig:ui_chatbot}
\end{figure}

Une pastille ronde apparaît en bas à droite de chaque page du portail Assiette Gala. D'un effleurement, le visiteur déplie une fenêtre de dialogue où il peut poser n'importe quelle question relative au menu, aux tarifs ou aux conditions de livraison. L'interface distingue visuellement les bulles du client (alignées à droite) de celles du robot (alignées à gauche). Un indicateur rotatif signale les secondes de réflexion de Gemini, et un bouton en croix referme le volet sans perdre l'historique.

\subsubsection{Conversation avec le chatbot}

\begin{figure}[H]
    \centering
    \includegraphics[width=0.5\textwidth]{ui_chatbot_conversation.png}
    \caption{Exemple de conversation avec le chatbot}
    \label{fig:ui_chatbot_conversation}
\end{figure}

\subsubsection{Recommandations personnalisées}

\begin{figure}[H]
    \centering
    \includegraphics[width=0.85\textwidth]{ui_recommendations.png}
    \caption{Interface de recommandations personnalisées}
    \label{fig:ui_recommendations}
\end{figure}

Sur la page d'accueil du portail, les habitués connectés découvrent une rangée horizontale de tuiles proposées par le moteur de suggestion. Chaque tuile présente la photographie d'un plat que le client n'a pas encore essayé mais qui ressemble à ses coups de c\oe{}ur antérieurs, accompagnée du prix et d'un bouton d'ajout express au panier.

% --------------------------------------------------
\section{Tests fonctionnels}

\begin{table}[H]
    \centering
    \caption{Tests fonctionnels --- Sprint 5}
    \label{tab:tests_sprint5}
    \renewcommand{\arraystretch}{1.3}
    \begin{tabular}{|c|p{3cm}|p{3.5cm}|p{2.5cm}|}
        \hline
        \textbf{N°} & \textbf{Cas de test} & \textbf{Comportement attendu} & \textbf{Résultat} \\
        \hline
        1 & Question sur un plat existant & Réponse avec informations du plat (prix, description, allergènes) & Réussi \\
        \hline
        2 & Question sur les catégories & Liste des catégories disponibles & Réussi \\
        \hline
        3 & Question sur les services traiteur & Description des services d'événements & Réussi \\
        \hline
        4 & Question hors périmètre & Réponse indiquant les limites du chatbot & Réussi \\
        \hline
        5 & Recommandations pour un client & Plats pertinents basés sur l'historique & Réussi \\
        \hline
        6 & Actualisation de la base de connaissances & Nouvelles données indexées et disponibles & Réussi \\
        \hline
        7 & Temps de réponse du chatbot & Réponse en moins de 5 secondes & Réussi \\
        \hline
        8 & Conversation multi-tours & Contexte maintenu entre les messages & Réussi \\
        \hline
    \end{tabular}
\end{table}

% --------------------------------------------------
\section{Conclusion}

Ce cinquième axe couronne le travail informatique par l'injection de matière neuronale dans le projet Assiette Gala Sfaxienne. Le greffon RAG permet au chatbot d'être incollable sur notre sélection traiteur, sans aucune dérive sur des sujets extérieurs. Grâce à cette boussole algorithmique accomplie de concert avec M. Elleuch, le service client devient omniprésent, confirmant l'aboutissement fonctionnel de ce système logiciel.
